The main tool used to construct 4--manifolds in later chapters is handle attachment.
We focus mainly on the definitions and results needed to meaningfully attach 1-- and 2--handles to a 4--manifold.

\begin{rmk}
	\label{rmk:corners}
	There is more than one way to define handle attachment, and we choose to do so in a way that feels more combinatorial in nature.
	The main concern with this approach is that the object resulting from handle attachment is a ``manifold with corners'' rather than a smooth manifold.
	There are arguments that the corners can be smoothed away in a canonical way such that a manifold obtained via handle attachment is smooth and unique up to diffeomorphism, but delving into such an argument at this point would be a distraction.
	A construction that does not require the smoothing of corners can be found in \cite{Kosi93}, but the machinery makes explicit handle attachment unnecessarily complicated.
\end{rmk}

There is a common tool in topology called the \emph{attaching map} that we use in this section and the next to build our machinery, so we take this opportunity to lay out the definition and notation.

\begin{defn}
  Let $X$ and $Y$ be topological spaces, $A\subset X$ a subspace, and $f:A\to Y$ a continuous map.
  We define a relation $\sim$ by putting $f(x)\sim x$ for every $x$ in $A$.
  Denote the quotient space $X\sqcup Y/\sim$ by $X\cup_f Y$.
  We call the map $f$ the \emph{attaching map}.  
  We say that $X$ is \emph{attached} or \emph{glued} to $Y$ over $A$.
  A space obtained through attachment is often called a \emph{adjunction space} or \emph{attachment space}.
\end{defn}

Throughout this section, we have $n=\lambda+\mu$, $M$ an $n$--manifold with boundary, and $H^\lambda=\D^\lambda\times\D^\mu$.
Attaching an $n$--dimensional $\lambda$--handle to $M$ is the process of joining $M$ to $H^\lambda$ along an embedding of $\pd\D^\lambda\times\D^\mu$ in $\pd M$. 

\begin{defn}[Handle]
	\label{def:handle}
	Let $\varphi:\pd\D^\lambda\times\D^\mu\to\pd M$ be an embedding, and consider it to be an attaching map between $M$ and $H^\lambda$.
	The space $M\cup_\varphi H^\lambda$ is a ``manifold with corners,'' and we smooth those corners immediately smooth in a canonical way as mentioned in Remark \ref{rmk:corners}.
	The attached space $H^\lambda$ is an \emph{$n$--dimensional $\lambda$--handle}, and $M\cup_\varphi H^\lambda$ is the result of an $n$--dimensional \emph{$\lambda$--handle attachment}.
\end{defn}

Note that the domain of $\varphi$ has the structure of a trivial $\mu$--disc bundle over $S^{\lambda-1}$ and, using the language of vector bundles, the image of $\varphi$ has the structure of a closed tubular neighbourhood of $f_0(S^{\lambda-1})=\varphi\comp z(S^{\lambda-1})$.
Our uniqueness theorems for tubular neighbourhoods tell us that the closed tubular neighbourhood
\[
	f:\D_{\pd M}f_0(S^{\lambda-1})\to\overline{\nu}_{\pd M}\,f_0(S^{\lambda-1})
\]
of $f_0(S^{\lambda-1})$ is unique up to ambient isotopy of the embedding $f_0$.
Thus, if we have an embedding $f_0:S^{\lambda-1}\to \pd M$ (i.e.\ a knot) with trivial disc bundle, then a handle attachment can be defined using $f_0$ and an explicit embedding 
\[
	f: S^{\lambda-1}\times\D^\mu\to \D_{\pd M}\,f_0(S^{\lambda-1})
\]
for a closed tubular neighbourhood $\overline{\nu}_{\pd M}\,f_0(S^{\lambda-1})$.

In other words, the characteristics of handle attachment that fully describe the smooth $n$--manifold $M\cup_\varphi H^\lambda$ up to diffeomorphism are:
\begin{enumerate}
	\item The isotopy class of an embedding $f_0:S^{\lambda-1}\to\pd M$ with trivial normal disc bundle, and 
	\item the isotopy class of an identification of $S^{\lambda-1}\times\D^\mu$ with $\overline{\nu}_{\pd M}\,f_0(S^{\lambda-1})$.
\end{enumerate}

There is some specific language that is useful to the description of handle attachment.

\begin{defn}
	Let $M\cup_\varphi H^\lambda$ be an $n$--manifold with a $\lambda$--handle attached.
	The embedding $f_0:S^{\lambda-1}\to\pd M$ is called the \emph{attaching map}, its image $f_0(S^{\lambda-1})$ is the \emph{attaching sphere}, and its tubular neighbourhood is the \emph{attaching neighbourhood}.
	The embedding $f: S^{\lambda-1}\times \D^\mu\to \overline{\nu}_{\pd M}\,f_0(S^{\lambda-1})$ is called a \emph{normal framing} or just \emph{framing}.
	Inside of the $\lambda$--handle $H^\lambda=\D^\lambda\times\D^\mu$, the disc $\D^\lambda\times\{\vec{0}\}$ is the \emph{core}, and the disc $\{\vec{0}\}\times\D^\mu$ is the \emph{cocore}.
	The boundary circle $\{\vec{0}\}\times S^{\mu-1}$ of the cocore is the \emph{belt sphere}.
	The integer $\lambda$ is the \emph{index} of the handle.
\end{defn}

It is desirable to describe a manifold entirely via handle attachments.
For this we define the general notion of a decomposition.

\begin{defn}
	\label{def:morsehandle}
	Let $W$ be a compact $n$--manifold.
	The boundary $\pd W$ will be taken to be a disjoint pair $M_0\sqcup M_1$, neither of which is necessarily nonempty.
	Describe $W$ as a chain of inclusions
	\[
	M_0 \times\I = W_0 \subset W_1 \subset W_2 \subset \cdots \subset W_{k-1} \subset W_k = W
	\]
	where $W_i$ is obtained from $W_{i-1}$ by attaching a handle and $M_0$ is identified with $M_0\times \{0\}$.
	This is called a \emph{handle decomposition} of the pair $(W,M_0)$.
\end{defn}

\begin{defn}
	A manifold $M$ with handle decomposition that consists of exactly one of 0--handle and $g$ 1--handles is called a \emph{handlebody} of \emph{genus} $g$.
	Let $V$ be a handlebody of genus $g$.
	A simple closed curve in $\pd V$ is called \emph{essential} if it is not homotopic to a point.
	A simple closed curve $J$, essential in $\pd V$, that bounds a 2--disc in $V$ is called a \emph{meridian}.
	The properly embedded disc in $V$ that has boundary $J$ is called a \emph{meridinal disc}.
	
	The special case of the oriented genus 1 handlebody is called a \emph{solid torus}.
	More generally, any space that is homeomorphic to $S^1\times\DD$ is called a \emph{solid torus}.
	A simple closed curve $J$ in the boundary of a solid torus that intersects a meridian at a single point is called a \emph{longitude}.
	Note that a longitude is essential in a solid torus and that there are infinitely many isotopy classes of longitudes, whereas there's exactly one isotopy class of meridians.
\end{defn}

An essential role of the handlebody is in the definition of a Heegaard splitting of a closed, orientable 3--manifold.

\begin{defn}	
	Let $U$ and $V$ be 3--dimensional handlebodies of genus $g$ and let $f:\pd U\to \pd V$ be an orientation reversing diffeomorphism.
	The adjunction $M=U\cup_f V$ is a \emph{Heegaard splitting} of $M$, the shared boundary $H = \pd U = \pd V$ is the \emph{Heegaard surface} of the splitting, and the shared genus of $U$ and $V$ is the \emph{genus} of the splitting as well.
	We also denote a Heegaard splitting by the pair $(M,H)$.
	
	We use the notion of equivalence between splittings from \cite{SchlWald}.
	A pair of splittings $(M,H)$ and $(M,H')$ are \emph{equivalent} if there is a homeomorphism $h:M\to M$ such that $h$ is isotopic to $\ident{M}$ and $\restr{h}{H}$ is an orientation preserving homeomorphism $H\to H'$.
\end{defn}

\begin{figure}
		\centering
		\caption{Genus 1 Heegaard splitting of $S^3$}
		\includegraphics[height=3cm]{figures/dummy.jpg}
		\label{fig:genus1split}
\end{figure}

\begin{ex}
	The 3--sphere $S^3$ has a number of standard Heegaard splittings.
	The first is the genus 0 splitting, usually realized by considering $S^3$ as the set of unit vectors in $\R^4$.
	Take the Heegaard surface to be the intersection of $S^3$ with the $xyz$--hyperplane in $\R^4$.
	This yields a copy of $S^2$ that separates $S^3$ into two connected components.
	This splitting is written $(S^3,S^2)$
	
	The second is the genus 1 splitting, which is visualized using the realization of $S^3$ as the one--point compactification of $\R^3$.
	Take solid tori $U$ and $V$ and identify $\pd U$ with $\pd V$ by the homeomorphism that swaps a meridian with a longitude.
	The adjunction that results is $S^3$, and the pair of solid tori are displayed in Figure \ref{fig:genus1split}.
	This splitting is written $(S^3,T^2)$.
	Note that this description of $S^3$ may also be obtained by examining the boundary of a 4--dimensional 2--handle $\pd (\DD\times\DD)$.
\end{ex}

The standard genus 1 splitting is used in conjunction with the connect sum operation in the following essential theorems.

\begin{defn}
	Let $M$ and $X$ be closed, oriented $n$--manifolds.
	The manifold obtained by identifying the boundary $(n-1)$--spheres of $M\setminus\B^3$ and $X\setminus\B^3$ in an orientation reversing way is called the \emph{connected sum} of $M$ and $X$ and is denoted $M\# X$.
\end{defn}

\begin{theorem}
	Connected sum is a well defined operation up to homeomorphism on the space of closed, oriented $n$--manifolds with identity element $S^n$.
\end{theorem}
	
\begin{defn}
	Let $(M,H)=U\cup_f V$ be a Heegaard splitting of $M$.
	The connected sum $(M,H)\#(S^3,T^2)$ is called an \emph{elementary stabilization} of $M$, and itself is a splitting $(M,H\#T^2)$.
	A Heegaard splitting $(M,H)$ is called a \emph{stabilization} of another splitting $(M,H')$ if it obtained from $(M,H')$ via a finite number of elementary stabilizations.
\end{defn}

Consider the meridians of the solid tori in the standard genus 1 splitting of $S^3$.
They each bound a disc in their respective handlebody, and they intersect in exactly one point.
We would expect to be able to find such curves in any 3--manifold obtained as a stabilization.
	
\begin{defn}	
	Let $(M,H)=U\cup_f V$ be a Heegaard splitting of genus $g$ and $\alpha$, $\beta$ a pair of simple, closed, essential curves in $H$.
	Let $\alpha$ be a meridian of $U$ and $\beta$ be a meridian of $V$ with associated meridinal discs $D_\alpha$, $D_\beta$.
	If $\alpha$ and $\beta$ intersect exactly once, then the pair $(\alpha,\beta)$ is a \emph{meridinal pair} or \emph{destabilizing pair} of the splitting.
\end{defn}

To see why $(\alpha,\beta)$ would be called a destabilizing pair, remove a tubular neighbourhood of $D_\alpha$ from $U$ and add it to $V$ as a 2--handle along a tubular neighbourhood of $\alpha$ in $H$.
In the proof of Waldhausen's Theorem it is shown that these altered spaces $U'$ and $V'$ are handlebodies of genus $g-1$ with $H'= \pd U' = \pd V'$, and $(M,H)$ is a stabilization $(M,H')\#(S^3,T^2)$.

\begin{theorem}[Waldhausen's theorem, \cite{SchlWald}]
	If $(S^3,H)$ and $(S^3,H')$ are Heegaard splittings of the same genus, then they are equivalent.
\end{theorem}

In the proof of Waldhausen's theorem, destabilizing pairs are identified and reduced.
We don't need the full power of Waldhausen's theorem, just a weaker version to use in Chapter \ref{cha:cobordisms}.

\begin{theorem}
	\label{thm:lilwald}
	Let $M$ be a closed, orientable 3--manifold with a Heegaard splitting $(M,H)$ of genus $g$.
	If there is a collection of $g$ disjoint destabilizing pairs in $H$, then $M$ is homeomorphic to $S^3$.
\end{theorem}

Heegaard splittings are a specific application of a specific type of handle decomposition.
General decompositions do not arise in this way.
They usually occur naturally by studying Morse functions as discussed in Chapter \ref{cha:cobordisms}.
More structure can be demanded out of a decomposition with handles of varying indices by isotoping attaching maps.

\begin{prop}[Proposition 4.2.7 of \cite{GompStip}]
	\label{prop:incrindex}
	Any handle decomposition of the compact pair $(W,M_0)$ can be modified so that handles are attached in order of increasing index.
	Handles of the same index can be attached in any order.	
\end{prop}

Because a $\lambda$--handle is also a $\mu$--handle of the same dimension, any handle decomposition actually defines a pair of dual decompositions.

\begin{defn}
	Let $W$ be a compact manifold with $\pd W = M_0 \sqcup M_1$ and handle decomposition
	\[
	M_0\times\I = W_0 \subset W_1 \subset W_2 \subset \cdots \subset W_{k-1} \subset W_k = W.
	\]
	By attaching a copy of $M_1 \times\I$ to $M_1$ and removing the collar $M_0\times\I$, we obtain a handle decomposition of the pair $(W,M_1)$.
	Each $\lambda$--handle in the decomposition of $(W,M_0)$ gives us a $\mu$--handle in the decomposition of $(W,M_1)$ that can be seen by reversing the roles of the core and cocore within the handle.
	The decomposition obtained is called the \emph{dual handle decomposition} of the pair $(W,M_1)$.
\end{defn}

Notice that if the handles of a decomposition are attached in order of increasing index, then the handles of the dual decomposition are as well.

There are two special cases of handle attachment.
The first is 1--handle attachment, which reduces to investigating the orientations of the attaching maps.
The second is 4--dimensional 2--handle attachment, which draws on concepts central to knot theory.

\begin{rmk}
	\label{rmk:1handle}
	Let $M$ be orientable and path--connected with $\pd M$ compact, connected, and nonempty.
	The attaching sphere of a 1--handle is $S^0=\pd \D^1$, which is a pair of points.
	There is a unique isotopy class of embeddings $f_0:S^0\to\pd M$.
	This means that $M\cup_\varphi H^1$ is determined entirely by the framing $f$.
	The normal disc bundle of $f_0(S^0)$ is a bundle over $S^0$, so it is vacuously trivial.
	Using the vector bundle structure of the tubular neighbourhood, we write an embedding of $S^0\times\D^{n-1}\to S^0\times\D^{n-1}$ as a pair of length--preserving linear transformations $\R^{n-1}\to\R^{n-1}$, i.e.\ elements of the orthogonal group $\orth{n-1}{\R}$, each restricted to act on one of the connected components of $S^0\times\D^{n-1}$.
	The determinant of an element of $\orth{n-1}{\R}$ is either $1$ or -$1$, and $\orth{n-1}{\R}$ has two path--connected components corresponding to these two cases.
	Every element of $\orth{n-1}{\R}$ is an embedding $\R^{n-1}\to\R^{n-1}$, so any path in $\orth{n-1}{\R}$ is an isotopy of its endpoints.
	It is then easy to see that are four isotopy classes of our trivialization, and those fall into two types.
	Either both transformations are are orientation preserving (reversing), or one is orientation preserving (reversing) and the other is orientation reversing (preserving).
	Under the first type of automorphism, $M\cup_\varphi H^1$ is a non-orientable manifold.
	Under the second, $M\cup_\varphi H^1$ is orientable.
\end{rmk}

Let $W$ be a 4--manifold with orientable boundary $M$.
To attach a 2--handle to $W$, we need an embedding $f_0:S^1\to M$ and a framing $S^1\times\D^2\to\overline{\nu}_{\pd M}f_0(S^1)$.
Recall that the space $S^1\times D^2$ is what we call a solid torus.
Once an embedding $f_0:S^1\to M$ is chosen, $W\cup_\varphi H^2$ is determined by the framing $f:S^1\times \D^2\to \overline{\nu}_{\pd M}\,f(S^1)$.
To investigate possible framings, we consider the \emph{mapping class group} of the solid torus.
The mapping class group of a space $X$ is the group of isotopy classes of automorphisms $X$, and it is denoted $mpg(X)$.
For $V$ a solid torus, $mpg(V)$ is very closely linked to the mapping class group of the torus $T^2 = S^1\times S^1=\pd V$, which is the special linear group $\textrm{SL}_2(\Z)$, the group of $2\times 2$ matrices with integer entries and determinant $\pm 1$.
\begin{lem}
	Let $V$ be a solid torus and let $f$ be an automorphism of the torus $\pd V$.
	Then $f$ extends to an automorphism of $V$ if and only if $f$ maps a meridian to a meridian.	
\end{lem}

\begin{lem}
	A pair of automorphisms $f,g$ of $V$ that agree on $\pd V$ and map a meridian to a meridian are isotopic.	
\end{lem}

\begin{theorem}
	The mapping class group of a solid torus $V$ is the subgroup of the mapping class group of $\pd V$ containing automorphisms that map meridians to meridians.
	This subgroup is isomorphic to $\Z$.	
\end{theorem}

\begin{proof}
	In \cite{Rolf76}, it is shown that the mapping class group of $T^2$ is $\textrm{SL}_2(\Z)$ by generating the group from a pair of basic automorphisms of $T^2$ called \emph{twists}, and the swapping map $(z,w)\to(w,z)$.
	The twists are defined as:
	\[
	\begin{array}{ccccc}
	h_L(e^{i\theta},e^{i\phi}) & = & (e^{i(\theta+\phi)},e^{i\phi}) & & \textrm{``longitudinal twist''} \\
	
	h_M(e^{i\theta},e^{i\phi}) & = & (e^{i\theta},e^{i(\theta+\phi)}) & & \textrm{``meridinal twist''}	
	\end{array}
	\]
	Fix the meridinal and longitudinal directions of the boundary of a solid torus $V$ to coincide with the meridinal and longitudinal directions of $V$.
	This would mean that if we are writing $V=S^1\times D^2$, $\{1\}\times S^1$ is a meridian and $S^1\times\{1\}$ is a longitude.
	Notice that the meridinal twist maps a meridian to a meridian, so it extends to an automorphism of $V$.
	Also, neither the longitudinal twist nor the swap map a meridian to a meridian, so neither extend to an automorphism of $V$.
	
	Let $F$ be an automorphism of $V$ with restriction $f$ to $\pd V$.
	Because $f$ is an automorphism of $\pd V$, it can be written as the product of twists and swaps.
	Because $f$ preserves meridians, it can be written entirely as a power of the meridinal twist.
	In the other direction, any automorphism that is the power of a meridinal twist clearly preserves meridians.
	Because these automorphisms are written as powers of $h_M$, there is a clear isomorphism from this subgroup to $\Z$.
\end{proof}

\begin{rmk}
	\label{rmk:2handle}
	The content of these results means that the attachment of a 2--handle to a 4--manifold can be described by an embedding $f_0:S^1\to M$ and an integer $k$ called the \emph{framing constant} that determines the framing automorphism as long as we know what a framing constant of $0$ means.
	To see this, first notice that $h_M$ can be extended to
	\[
	\begin{array}{cccc}
		H_M: & S^1\times\DD & \to & S^1\times\DD \\
		& (e^{i\theta}, re^{i\phi}) & \mapsto & (e^{i\theta},re^{i(\phi+\theta)})	
	\end{array}
	\]
	which represents a generator of $mpg(S^1\times\DD)$.
	Because $\nu_{\pd M}\,f(S^1)$ is a solid torus, a one-to-one correspondence
	$\psi$ between $mpg(S^1\times\R^2)$ and the isotopy classes of maps $S^1\times\DD\to\nu_{\pd M}\,f(S^1)$ is determined entirely by $\psi(\id)$.
	The image of $\psi(\id)$ is called the class of \emph{preferred framings}.
	Sometimes, there is a canonical class of preferred framings.
	One important case is when $M$ is $S^3$.
	
	Let $f_0:S^1\to S^3$ be an embedding.
	Let $V$ be a solid torus that is a closed tubular neighbourhood of the knot $f_0(S^1)$ in $S^3$.
	Then there is exactly one isotopy class $[J]$ of longitudes of $V$, unique up to ambient isotopy of $f_0$, such that any curve in $[J]$ bounds a disc in $S^3\setminus V$.
	Moreover, for any representative $J$ of $[J]$ there is a framing $f_J:S^1\times\DD\to V$ such that $f_J(S^1\times\{\vec{1}\})=J$, and the isotopy class of $[f_J]$ is in one--to--one correspondence with $[J]$.
	Our canonical choice for $\psi(\id)$ is $[f_J]$, and this generates the rest of the map, i.e.\ $\psi([H_M^k])=[f_J\comp H_M^k]$.
	The isotopy class of a framing $f:S^1\times\DD\to\overline{\nu}_{S^3}f_0(S^1)$ can then be classified entirely by an integer $k$ that describes the number of meridinal twists it takes to get to $f$ from the preferred one.
	We call $k$ the \emph{framing constant} for $f$.
\end{rmk}
